\documentclass[10pt]{article}

\usepackage[utf8]{inputenc}

\usepackage[T1]{fontenc}

\usepackage{amsmath}

\usepackage{amsfonts}

\usepackage{amssymb}

\usepackage[version=4]{mhchem}

\usepackage{stmaryrd}

\usepackage{bbold}

\usepackage{graphicx}

\usepackage[export]{adjustbox}

\graphicspath{ {./images/} }



\begin{document}

Изучение М.Ф. началось в 19 в. в связи с изучением эллиптич. функщий и предшествовало появлению общей теории автоморфных функций. В теории М. Ф. в качестве основных модулярных форм используются следующие ma-ряды:



$g_{2}(z)=60 \sum_{m_{1}, m_{2} \in \mathbb{Z}}^{*} \frac{1}{\left(m_{1}+m_{2} z\right)^{4}}=(2 \pi)^{4}\left\{\frac{1}{12}+20 \sum_{n=1}^{\infty} \frac{n^{3} h^{2 n}}{1-h^{2 n}}\right\}$,



$g_{3}(z)=140 \sum_{m_{1}, m_{2} \in \mathbb{Z}}^{*} \frac{1}{\left(m_{1}+m_{2} z\right)^{6}}=(2 \pi)^{6}\left\{\frac{1}{216}-\frac{7}{3} \sum_{n=1}^{\infty} \frac{n^{5} h^{2 n}}{1-h^{2 n}}\right\}$, $\Delta(z)=g_{2}^{3}(z)-27 g_{3}^{2}(z)$,



где $h=e^{i \pi z}$, а звездочка означает, что нулевая пара $\left(m_{1}, m_{2}\right)=(0,0)$ отбрасывается. Согласно терминологии K. Вейерштрасса (K. Weierstrass), это - от нос и тель н ы е и н в а р и а н ты, играющие важную роль в его теории эллиптич. функщий (см. Вейерштрасса эллиптические функ$ц u u)$, а $\Delta$ называется также д и с к р и м и на н т о м. С точки зрения теории автоморфных функщий, это - автоморфные формы соответственно веса 2, 3 и 6, ассоциированные с модулярной группой. Основная М. Ф. имеет вид:



\[

J(z)=\frac{g_{2}^{3}(z)}{\Delta(z)}=1+\frac{27 g_{3}^{2}(z)}{\Delta(z)}

\]



Функция $J(z)$ называется также а 6 с ол ю т ны м и н в а р и а н т о м. Она регулярна в верхней полуплоскости, а внутри фундаментальной области $G$ принимает каждое конечное значение, кроме 0 и 1, в точности один раз; кроме того, $J(-1 / 2+i \sqrt{3} / 2)=0, J(i)=1$.



В теории эллиптич. функций М. Ф. $J(z)$ играет важную роль, позволяя по заданным вейерштрассовым относительным инвариантам $a=g_{2}, b=g_{3}, a^{3}-27 b^{2} \neq 0$ определить периоды $2 \omega_{1}, 2 \omega_{3}$, а следовательно, и построить все эллиптич. функции Вейерштрасса. Если $\tau$ - единственное в фундаментальной области решение уравнения



\[

J(\tau)=a^{3} /\left(a^{3}-27 b^{2}\right)

\]



то при $a \neq 0, b \neq 0$ имеем $\omega_{1}^{2}=a / b, \omega_{3}=\omega_{1} \tau$; при $a=0$ имеем $\tau=-1 / 2+i \sqrt{3 / 2}$, и $\omega_{1}$ определяется из уравнения



\[

\omega_{1}^{6}=\frac{140}{b} \sum_{m_{1}, m_{2} \in \mathbb{Z}}^{*} \frac{1}{\left(m_{1}+m_{2} \tau\right)^{6}}, \quad \omega_{3}=\omega_{1} \tau

\]



при $b=0$ имеем $\tau=i$, и $\omega_{1}$ определяется из уравнения



\[

\omega_{1}^{4}=\frac{60}{a} \sum_{m_{1}, m_{2} \in \mathbb{Z}}^{*} \frac{1}{\left(m_{1}+m_{2} \tau\right)^{4}}, \quad \omega_{3}=\omega_{1} \tau .

\]



Для построения Якоби эллиптических функций удобнее, вместо $J(z)$, использовать функцию



\[

\lambda(z)=k^{2}(z)=1-\prod_{n=1}^{\infty}\left(\frac{1-h^{2 n-1}}{1+h^{2 n+1}}\right)^{8}, \quad h=e^{\pi i z},

\]



\begin{center}

\includegraphics[max width=\textwidth]{2023_07_08_57b6f77dfe0d6c5f7168g-1}

\end{center}



также называемую М. Ф. На самом деле $\lambda(z)$ является автоморфной функцией только относительно подгруппы $\Gamma_{2}$ модулярной группы $\Gamma$, причем к $\Gamma_{2}$ относятся все те преобразования вида (1), у к-рых (в качестве дополнительного условия) $a$ и $d$ - нечетные числа, $b$ и $c$ - четные. Фундаментальная область $G_{2}$ группы $\Gamma_{2}$ изображена на рис. 2; это - криволинейный четырехугольник $A B O C A$ с вершинами $A(\infty), B(-1), O, C(1)$, две стороны к-рого $\boldsymbol{A B}$ и



366 МОДУЛЯРНАЯ $C A$ - отрезки прямых соответственно $x=-1$ и $x=1$, а $B O$ и $O C$ - дуги окружностей соответственно $|z+1 / 2|=1 / 2$ и $|z-1 / 2|=1 / 2$. Участки границы слева от мнимой оси включаются в $G_{2}$, а $O C$ и $C A$ не включаются. Функция $\lambda(z)$ также регулярна в верхней полуплоскости $\operatorname{Im} z>0$. Внутри области $G_{2}$ она принимает каждое конечное значение, кроме 0 и 1 , в точности один раз; кроме того, $\lambda(\infty)=0$ и $\lambda(0)=1$.



Лит.: [1] Гурвиц А., Курант Р., Теория функций, пер. [с нем.], М., 1968; [2] А х и е зе р Н.И., Элементы теории эллиптических функций, 2 изд., М., $1970 . \quad E$. Д. Соломенцев. MOMEHT случайной величины $X-$ математическое ожидание ее степеней; мо м е н т порядка $n=1,2,3, \ldots$ для непрерывно распределенной случайной величины с плотностью $p(x)$ равен



\[

\mathrm{E} X^{n} \equiv\left\langle X^{n}\right\rangle=\int_{-\infty}^{\infty} x^{n} p(x) d x

\]



Для дискретной случайной величины, принимающей значения $\left\{x_{k}\right\}$ с вероятностями $\left\{p_{k}\right\}, n$-й М. равен



\[

\mathrm{E} X^{n}=\sum_{k} x_{k}^{n} p_{k}

\]



M. 1-го порядка $\mathrm{E} X$ - математич. ожидание. Величина



\[

\mathrm{E}(X-\mathrm{E} X)^{n}

\]



называется центральным моментом порядка $n$, центральный М. 2-го порядка называется дисперсией.



В случае конечного или бесконечного семейства случайных величин $\left\{X_{t}, t \in T\right\}$, где $T-$ нек-рое множество, помечающее эти величины, с мешан ны е моменты (м ульт и м о м е н т ы) этого семейства определяются формулой



\[

M_{n_{1}, \ldots, n_{k}}\left(t_{1}, \ldots, t_{k}\right)=E X_{t_{1}}^{n_{1}} \ldots X_{t_{k}}^{n_{k}}

\]



где $\left\{t_{1}, \ldots, t_{k}\right\}-$ произвольный набор попарно различных точек $T$, а $\left\{n_{1}, \ldots, n_{k}\right\}-$ целочисленный мультииндекс; математич. ожидание вычисляется по совместному распределению вероятностей значений случайных величин $X_{t_{1}}, \ldots, X_{t_{k}}$



Р. А. Минлос.



МОМЕНТ ИНЕРЦИИ - величина, характеризующая распределение масс в теле и являющаяся наряду с массой мерой инертности тела при непоступательном движении. В механике различают М. и. осевые и центробежные. Осевым моментом инерции тела относительно оси $z$ называется величина, определяемая равенством:



\[

I_{z}=\sum m_{i} h_{i}^{2} \text { или } I_{z}=\int_{V} \rho h^{2} d V

\]



где $m_{i}$ - массы точек тела, $h_{i}-$ их расстояния от оси $z$, а $\rho-$ массовая плотность, $V$ - объем тела. Величина $I_{z}$ является мерой инертности тела при его вращении вокруг оси. Осевой М. и. можно также выразить через линейную величину $r$, называемую р а д и у с ом и н е р ц и и, по формуле



\[

I_{z}=M r^{2}

\]



где $M$ - масса тела. Размерность М. и. $-L^{2} M$.



Центробежными моментами инерции относительно системы прямоугольных осей $x, y, z$, проведенных в точке $O$, называются величины, определяемые равенствами:



\[

\left.\begin{array}{rl}

I_{x y} & =\sum m_{i} x_{i} y_{i}, \\

I_{y z} & =\sum m_{i} y_{i} z_{i}, \\

I_{z x} & =\sum m_{i} z_{i} x_{i}

\end{array}\right\}

\]



или же соответствующими объемными интегралами. Эти величины являются характеристиками динамич. неуравновешенности тел. Напр., при вращении тела вокруг оси $z$ от значений $I_{x z}$ и $I_{y z}$ зависят силы давления на подшипники, в к-рых закреплена ось.



M. и. относительно параллельных осей $z$ и $z^{\prime}$ связаны соотношением:



\[

I_{z}=I_{z^{\prime}}-M d^{2}

\]





\end{document}